% !TEX root = ../jlt.tex
\section{From \Cpp{} to Jlt}
\label{sec:posture}

Jlt is standard \Cpp{23} under three rules, not a dialect and not a library; the compiler enforces some of the rules, and we say which.
Each rule has a listing of a few lines, verbatim from the repository's test up to formatting and the definitions a listing omits or the test imports; the LLVM IR of Appendix~\ref{app:fjlt:ssa} is Clang~23's output, its flags in its caption, with names shortened.

\paragraph{Rule 1, Substitution: values and pure functions.}
\phantomsection\label{def:rule1}
Every type a Jlt program passes around is \cppinline{std::semiregular}: default constructible, copyable, and nothing else is assumed.
Every function is \cppinline{constexpr} and depends on its arguments alone.\footnote{A \cppinline{constexpr} function may still change its own locals: constant evaluation permits modifying only objects whose lifetime began within it, so no mutation can leak, which is the encapsulated state of \cppinline{runST}~\cite{launchbury1994lazy}.
At runtime the linter of Section~\ref{sec:limits} holds the body to the same condition by closing each channel a mutation could leak through: parameters, results, the object, statics, globals, \cppinline{volatile}, boxes, and escaping closures that capture by reference.}
What the rule conserves is what has come to be called Leibniz's law, that equals cannot be told apart by any operation, at either binding time, whose converse is the principle Leibniz himself stated~\cite{leibniz1686discours,leibniz1686discourse}: a call may be replaced by its value, $e[f(x)] = e[v]$ whenever $f(x)$ evaluates to $v$.
In a programming language the law is referential transparency~\cite{strachey2000fundamental}, and a language where $a = b$ but $f(a) \neq f(b)$ for some $f$ has an equality one cannot reason with.

\begin{listing}[htbp]
\begin{minipage}[t]{0.48\linewidth}
\begin{cpp}
// Identity arrows:
template <typename A>
constexpr auto id() { return Identity<A>{}; }
constexpr auto id_b = id<bool>();
constexpr auto id_u = id<size_t>();
static_assert(
id_b(true) == true && id_u(42) == 42 &&
is_monotone(id_u) && is_monotone(id_b) &&
is_involution(id_u) && is_involution(id_b) &&
is_identity(id_u >> id_u) &&
(id_u >> id_u)(42) == 42 &&
is_monotone(id_u >> id_u) &&
is_involution(id_u >> id_u));
\end{cpp}
\end{minipage}\hfill
\begin{minipage}[t]{0.48\linewidth}
\begin{cpp}
// Constant arrows:
constexpr auto nil = constant<bool>(0uz);
constexpr auto one = constant<bool>(1uz);
constexpr auto ans = constant<bool>(42uz);
constexpr auto aye = constant<size_t>(true);
constexpr auto nay = constant<size_t>(false);
static_assert(
(ans >> id_u)(true) == ans(false) &&
is_constant(ans) && is_monotone(ans) &&
is_constant(aye) && is_antitone(aye) &&
!is_monic(ans) && !is_epic(ans) &&
!is_monic(aye >> id_b) &&
!is_epic(nay >> id_b) &&
is_constant(nil >> nay >> ans));
\end{cpp}
\end{minipage}\hfill
\begin{minipage}[t]{0.48\linewidth}
\begin{cpp}
// Involutions, pointwise:
static_assert(
(!aye)(0) == false &&
is_involution(!id_b) &&
is_constant(!aye) &&

(~nil)(true) + 1 == 0 &&
is_involution(~id_u) &&
is_constant(~nil) &&

(-one)(true) == (~nil)(true) &&
is_involution(-id_u) &&
is_constant(-nil) &&

is_iso(-id_u >> ~id_u) &&
!is_involution(-id_u >> ~id_u));
\end{cpp}
\end{minipage}\hfill
\begin{minipage}[t]{0.48\linewidth}
\begin{cpp}
// Products and coproducts:
constexpr auto fst = Π<1, pair<bool, size_t>>{};
constexpr auto sec = Pairing{id_b, nil};
constexpr auto inl = ι<1, bool, size_t>{};
constexpr auto ret = Copairing{id_b, nay};
static_assert(
is_epic(fst) && !is_monic(fst) &&
is_monic(sec) && !is_epic(sec) &&
(sec >> fst)(true) && !(sec >> fst)(false) &&
is_monic(inl) && !is_epic(inl) &&
is_epic(ret) && !is_monic(ret) &&
(inl >> ret)(true) && !(inl >> ret)(false));
\end{cpp}
\end{minipage}
\caption{\hyperref[def:rule1]{Rule~1}: Substitution.
Equals cannot be told apart, so a function is its graph, and functions are the arrows of a category, with composition \cppinline{>>} their only structure.
Top left: the identities.
Top right: the points, constant arrows, monotone and antitone at once and neither monic nor epic.
Bottom left: the involutions, each the C++ operator post-composed, so \cppinline{!f} is $\neg \circ f$; two compose to an isomorphism and to no involution.
Bottom right: a projection and its section, an injection and its retraction, the arrows that are epic or monic but not both.
A law is asked on the arrow's value, at compile time: \cppinline{is_monic(f)} reads what the concept \cppinline{IsMonicArrow} reads, and a composite's law is derived from its legs.}
\label{lst:rule1}
\end{listing}

In Listing~\ref{lst:rule1}, \cppinline{const T} and \cppinline{const T&} are one parameter, the second an optimisation the translation erases, since nothing changes the value the reference denotes.
An arrow is a function made a value, and its laws are equations between values, so \cppinline{==} is where Leibniz's law meets the program: the kernel asks it of no type, the listings of every carrier a law is stated on.
The properties a law singles out, mono, epi and iso by cancellation~\cite{pierce1991basic}, involution, the variances, are certificates declared or derived on the arrow, whose type names carrier and operation, so a certificate is about the signature.
Two consequences keep that exact: a function returns a value, never a reference, and every member function, a call operator included, is \cppinline{const}.
The grammar of Appendix~\ref{app:fjlt:syntax} has neither a reference type nor a non-\cppinline{const} member, and the linter of Section~\ref{sec:limits} holds \Cpp{} to the same.
The compiler enforces the rule in part: a \cppinline{constexpr} function that reaches for state fails to evaluate at compile time, but nothing stops a caller from only ever calling it at runtime.

\paragraph{Rule 2, Abstraction: a type varies behaviour only through its concepts and traits.}
\phantomsection\label{def:rule2}
Generic code may do with a \cppinline{T} what \cppinline{T}'s concepts and traits say it can, and nothing else.
That is the premise of Wadler's free theorems~\cite{wadler1989theorems}, with concepts and traits in the place of type classes~\cite{wadler1989adhoc}, and what it conserves is the body's independence from the carrier: \cppinline{f<T>} is one generic term applied to \cppinline{T}'s dictionary, so it respects every relation the dictionary respects (Corollary~\ref{cor:parametricity}).

\begin{listing}[htbp]
\begin{minipage}[t]{0.47\linewidth}
\begin{cpp}
using namespace std;

template <typename T, typename Op>
constexpr T unit(Op) = delete;
template <typename T, typename Op>
constexpr bool associative(Op) = delete;

template <unsigned_integral T>
constexpr T unit(plus<T>) { return 0; }
template <unsigned_integral T>
constexpr bool associative(plus<T>) {
  return true;
}
template <typename T, typename Op>
concept HasUnit = requires(Op op) {
  { unit<T>(op) } -> same_as<T>;
};
template <typename T, typename Op>
  requires HasUnit<T, Op> && equality_comparable<T>
constexpr bool is_unit(const T& a, Op op) {
  return a == unit<T>(op);
}

template <typename T, typename Op>
concept IsMonoid =
    semiregular<T> &&
    requires(const T a, const T b) {
      { Op{}(a, b) } -> same_as<T>;
      { unit<T>(Op{}) } -> same_as<T>;
      { is_unit(a, Op{}) } -> same_as<bool>;
    } && associative<T>(Op{});
\end{cpp}
\end{minipage}\hfill
\begin{minipage}[t]{0.5\linewidth}
\begin{cpp}
template <typename T, typename Op = plus<T>>
  requires IsMonoid<T, Op>
constexpr T concat() { return unit<T>(Op{}); }
template <typename T, typename Op = plus<T>>
  requires IsMonoid<T, Op>
constexpr T concat(T w) { return w; }
template <typename T, typename Op = plus<T>>
  requires IsMonoid<T, Op>
constexpr T concat(T v, T w) { return Op{}(v, w); }

static_assert(
concat<size_t>() == 0 &&
concat<size_t>(42uz) == 42 &&
concat<size_t>(40uz, 2uz) == 42 &&
concat<size_t, multiplies<size_t>>(6uz, 7uz)
    == 42 &&
concat<string>() == "" &&
concat<string>("ab"s, "c"s) == "abc" &&
concat<bool, logical_or<bool>>(false, true) &&
concat<bool, logical_and<bool>>(true, true) &&
concat<Sections<size_t>, compose>()(42uz)
    == 42 &&
concat<Sections<size_t>, compose>(
    adding(40uz), adding(2uz))(0uz) == 42);
\end{cpp}
\end{minipage}
\caption{\hyperref[def:rule2]{Rule~2}: Abstraction.
Left: a monoid is a set which is closed under an associative binary operation and equipped with a unit element.
Its elements need not be comparable, since equality is not decidable on every carrier, on functions for one.
The unit is identified instead: \cppinline{is_unit}, by default equality with the unit, which a carrier may supply otherwise.
The monoid shape is certified for the family of \cppinline{std::unsigned_integral} under \cppinline{(0, +)}.
Right: \cppinline{concat} at the lengths that need no recursion, none, one and two.
The monoid signatures are: \cppinline{size_t} under \cppinline{(0, +), (1, *)}, \cppinline{bool} under \cppinline{(false, ||), (true, &&)}, \cppinline{string} under \cppinline{("", +)}, and \cppinline{Sections}, the sections of $+$, under \cppinline{(id, >>)}, with \cppinline{adding} the curried $+$.
}
\label{lst:rule2}
\end{listing}

\begin{figure}[tbp]
\centering
\begin{tikzcd}[column sep=2em, row sep=1.5em]
  \llap{$\mathbf{1}$:\quad}
  &[1.6em] & & & & & & & & \mathbf{1} \arrow[d, "n"] \\
  \llap{$X$:\quad} \ast
  & 0 \arrow[l, "\mathit{prior}"'] \arrow[d, "f"]
  & & 1 \arrow[ll, "\mathit{prior}"'] \arrow[d, "f"]
  & & \cdots \arrow[ll, "\mathit{prior}"']
  & & n-1 \arrow[ll, "\mathit{prior}"'] \arrow[d, "f"]
  & & n \arrow[ll, "\mathit{prior}"'] \arrow[rr, dashed, "\mathit{succ}^*"]
 % & & \cdots \arrow[rr, dashed, "\mathit{succ}"]
  & & \mathbb{N} \\
  \llap{$S$:\quad}
  & f(0) \arrow[d]
  & & f(1) \arrow[d]
  & & \cdots
  & & f(n-1) \arrow[d]
  & \\
  \llap{$M$:\quad}
  & \bullet
  & & \circ \arrow[ll, "h"']
  & & \cdots \arrow[ll, "h"']
  & & \circ \arrow[ll, "h"']
  & & e \arrow[ll, "h"']
\end{tikzcd}
\caption{Where each structure lives.
Seed: after Lambek~\cite{lambek1968fixpoint}, the structure map $[0, \mathit{succ}] : \mathbf{1} + X \to X$ of the naturals is an isomorphism.
Its inverse $\mathit{prior} : X \to \mathbf{1} + X$ steps down and always ends.
Peano's $\mathit{succ}$~\cite{peano1889arithmetices}, dashed, never does.
The user supplies the point $n : \mathbf{1} \to X$ and $\mathit{prior}$.
The chain is what $\mathit{prior}$ reaches from $n$: the $n$ elements below it, not $n$ itself.
So $n$ names the chain's cardinality, and $0$ the empty chain.
Order: after Dedekind~\cite{dedekind1888zahlen}, a chain is a set closed under a map into itself.
The chain of an element is its forward orbit, here the chain of $\mathit{prior}\,n$.
The carrier needs neither an order nor an origin.
Reachability under $\mathit{prior}$ orders the chain, and its bottom is where $\mathit{prior}$ has nothing.
Algebra: $f$ lands in $S$, and the fold in $M$.
In the simplest case $S = M$ is a monoid.
A monoid $S$ may also act on $M$, and a semiring on a semimodule, which adds distributivity.
Fold: a list has two folds~\cite{gibbons1996third}.
The one along the walk, $h : M \times S \to M$, runs from $e$ at the seed's end through the intermediate values, $\circ$, to the result, $\bullet$.
It needs no list: it is a stream and its iterator~\cite{stepanov2009elements}.
With $\mathit{succ}$ as well, $\mathit{succ} >> \mathit{prior} = \eta$, the chain is a zipper~\cite{huet1997zipper}.
With offsets it is an array~\cite[\S 2.2.2]{knuth1997art1}.
These are the standard's forward, bidirectional and random-access iterators~\cite{cpp23standard}.
}
\label{fig:structures}
\end{figure}

\cppinline{concat} uses only what its constraint names, so it is a function of the monoid's record.
A \Cpp{} compiler checks a template's body at instantiation, against the type supplied, and not at definition, against the constraint, which is why this is a rule here and not a theorem.

\paragraph{Rule 3, Descent: recursion is structural.}
\phantomsection\label{def:rule3}
A Jlt program recurses on a structure whose residue, what remains to be folded, gets smaller at every step and so reaches its end.
The step is Figure~\ref{fig:structures}'s $\mathit{prior}$, which steps down and has nothing at the bottom, or another step down the same order, as halving is (Listing~\ref{lst:laws}), on an order with no infinite descent.
What the rule conserves is that inductive argument: structural induction~\cite{burstall1969proving}, a language all of whose programs end by it being, in a stricter form, Turner's total functional programming~\cite{turner2004total}, and the fold's universal property its equation~\cite{meijer1991functional,hutton1999tutorial}, $h = \mathit{fold}\,a$ exactly when $h \circ \mathit{in} = a \circ F\,h$.\footnote{$F$ is one layer of the data type with a hole for the rest, for a list of words $F\,X = 1 + \mathit{Word} \times X$; $\mathit{in}$ builds one more layer, the constructors $\mathit{nil}$ and $\mathit{cons}$; $a$, the algebra, says what to do with one layer once the rest is folded; and $F\,h$ applies $h$ to the rest only.
For \cppinline{concat} the equation is its definition, $h(\mathit{nil}) = \mathit{unit}$ and $h(\mathit{cons}(w, r)) = w \oplus h(r)$, and \emph{exactly when} says that \cppinline{concat} is the only function that satisfies both.}
Meijer et al.\ take their recursion operators in CPO, where one list type holds the finite, the infinite and the partial lists, and a fold over it need not end; Jlt takes them for data only, the initial algebras in Set, where a list in a type is finite as an initial algebra and the chain in a value is finite by the descent, and codata is out of scope: a coalgebra is admitted as a generator bounded by a count, never as a value.
The structure lives in a \emph{type} (a typed term, an array of known length), or in a \emph{value}.
In both cases, an iteration that can be mapped onto the chain of Figure~\ref{fig:structures} contains no loop and terminates, whether the chain is held as one count or as the elements of an array.
In the terms of Meijer et al., \cppinline{concat} over a list is a catamorphism.
The chain is an anamorphism, $\mathit{prior}$ unfolded from the seed and delimited where it has nothing.
\cppinline{concat} over the chain is a hylomorphism, and the chain is never built as a list, only as its fold (Appendix~\ref{app:fjlt:translation}).

\begin{listing}[htbp]
\begin{minipage}[t]{0.49\linewidth}
\begin{cpp}
// the counts, a finite chain
static_assert(IsFinite<size_t> &&
    IsTotalOrder<size_t, less_equal<>>);
constexpr auto ns =                   // 9, …, 0
    10uz >> chain<size_t>();
static_assert(
    ranges::input_range<decltype(ns)> &&
    ranges::distance(ns) == 10 &&
    ranges::equal(ns, views::iota(0uz, 10uz)
                      | views::reverse) &&
    ranges::empty(0uz >> chain<size_t>()) &&
    ranges::equal(ns >> id<size_t>(), ns));
\end{cpp}
\end{minipage}\hfill
\begin{minipage}[t]{0.49\linewidth}
\begin{cpp}
// Collatz's step, until 1
constexpr auto prior = until(equals(1uz),
    arrow<size_t>([](size_t n) {
      return n % 2 ? 3 * n + 1 : n / 2;
    }));
constexpr IsBoundedRange auto cs = // 500, 250, …, 1
    1'000uz >> chain(prior, 200uz);
static_assert(
    bound(cs) == 200 &&
    ranges::distance(cs) == 111 &&
    ranges::max(cs) == 9232);
\end{cpp}
\end{minipage}\hfill
\begin{minipage}[t]{0.49\linewidth}
\begin{cpp}
// the fold on a monoid: the library's concat
template <IsBoundedRange R,
          typename T = ranges::range_value_t<R>,
          typename Op = plus<T>>
  requires IsMonoid<T, Op>
constexpr T concat(R r, Op op = {}) {
  return ranges::fold_left(r, unit<T>(op), op);
}
static_assert(
    concat(ns) == 45 &&
    concat(ns, multiplies<size_t>{}) == 0 &&
    concat(ns >> constant<size_t>(7uz)) == 70 &&
    concat(cs) == 102512);
\end{cpp}
\end{minipage}\hfill
\begin{minipage}[t]{0.49\linewidth}
\begin{cpp}
// the fold on an action: matrices on vectors
using V = array<size_t, 2>;
struct M { V top, bottom;           // row by row
  friend constexpr V operator*(V v, M a) {  // v·A
    return {v[0] * a.top[0] + v[1] * a.bottom[0],
            v[0] * a.top[1] + v[1] * a.bottom[1]};}
  /* M * M, M + M, n * M, == */ };
constexpr auto act = [](V v, M a) { return v * a; };
constexpr M fib{{1, 1}, {1, 0}};
constexpr auto ms = ns >> constant<size_t>(fib);
// F(11), F(10)
constexpr auto fs_10 = ranges::fold_left(
  ms, V{1, 0}, act);
static_assert(fs_10 == V{89, 55}); 
\end{cpp}
\end{minipage}
\caption{\hyperref[def:rule3]{Rule~3}: Descent, Figure~\ref{fig:structures} in code.
Top. Left: counting down, a finite chain, whose type supplies $\mathit{prior}$.
Right: a chaotic path, Collatz's step from $1000$.
It is finite by fact, as far as anyone knows: that every natural reaches $1$ is an open problem~\cite{lagarias1985collatz}.
So its chain declares a budget, and is bounded by it.
Bottom. Left: a closed operation, \cppinline{concat} on any monoid, generically, Listing~\ref{lst:rule2}'s at a range.
Right: a general fold is a heterogeneous action, here matrices on vectors, Fibonacci's.
In the standard's ranges library~\cite{cpp23standard}, Figure~\ref{fig:structures}'s chain is an \cppinline{input_range}, and \cppinline{IsBoundedRange} one with a count known before the fold starts.
The arrow $f$ along it is \cppinline{views::transform}, and the fold with the action $h$ is \cppinline{ranges::fold_left}.}
\label{lst:rule3}
\end{listing}

What the rule buys is termination and, with Rule~2, the System~F reading, data as their own folds after B\"ohm and Berarducci~\cite{bohm1985automatic} (Appendix~\ref{app:fjlt:translation}).
Like System~F, strongly normalising and so unable to express every computable function~\cite{girard1989proofs}, Rule~3 trades expressiveness for termination, since the halting programs are not a decidable set~\cite{turing1936computable}: a bounded search ends without always deciding its question (Section~\ref{sec:binding}).

\begin{table}[ht]
\centering
\small
\begin{tabularx}{\linewidth}{@{}l>{\raggedright\arraybackslash}X>{\raggedright\arraybackslash}X>{\raggedright\arraybackslash}X>{\raggedright\arraybackslash}X@{}}
\toprule
rule & archetype & keeps out of \Cpp{} & buys & costs \\
\midrule
1 Substitution & \texttt{id}, the identity arrow & state, pointers, exceptions & one program, two binding times & no state beyond a call's own locals, no I/O \\
2 Abstraction & \texttt{concat}, a monoid's fold at a length in its type & a body using what its constraints do not name & the System~F reading & no ad hoc use of a type \\
3 Descent & \texttt{concat} over a chain, a fold over structure in a value & unbounded loops, general recursion & termination & a declared budget, for a chain finite only by fact \\
\bottomrule
\end{tabularx}
\caption{The posture at a glance, each rule named for what it conserves, with what it keeps out of \Cpp{}: the filtration from \Cpp{} to Jlt.
At build time, Rule~1 is held by the compiler in part and by the repository's linter for its syntax; Rule~2 by the linter at the level of names and by an auditing reader for the rest; Rule~3 by the compiler at compile time and by the linter for loops and recursion.
Section~\ref{sec:limits} says how.}
\label{tab:posture}
\end{table}

